\bigskip


Une exposition d'art contemporain a lieu dans une salle en forme de pavé droit de largeur 6 m, de longueur 8 m et de hauteur 4 m. 

Elle est représentée par le parallélépipède rectangle OBCDEFGH où
OB = 6 m, OD = 8 m et OE = 4 m.

On utilise le repère orthonormé \Oijk{} tel que $\vect{\imath} = \dfrac16\vect{\text{OB}}, \vect{\jmath} = \dfrac18\vect{\text{OD}}$ et $\vect{k} =\dfrac14\vect{\text{OE}}$.

\begin{center}
\psset{unit=1cm}
\begin{pspicture}(7.2,5.3)

\pspolygon(0.2,1.2)(4.5,0.2)(4.5,2.7)(0.2,3.7)
\psline(4.5,0.2)(6.9,1.2)(6.9,3.7)(4.5,2.7)
\psline(6.9,3.7)(2.6,4.7)(0.2,3.7)
\psline(0.2,2.45)(1.8,3.33)(1.4,4.2)
\psline[linestyle=dashed,linewidth=1.5pt](0.2,2.45)(1.4,4.2)
\psline[linestyle=dotted,linewidth=1.5pt](0.2,1.2)(2.6,2.2)(2.6,4.7)
\psline[linestyle=dotted,linewidth=1.5pt](2.6,2.2)(6.9,1.2)
\psline{->}(2.6,2.2)(2.2,2.02)
\psline{->}(2.6,2.2)(3.4,2.03)
\psline{->}(2.6,2.2)(2.6,2.82)
\psdot(2.7,1.5)
\uput[l](0.2,2.45){A} \uput[dl](0.2,1.2){B} \uput[d](4.5,0.2){C} \uput[dr](6.9,1.2){D}
\uput[u](2.6,4.7){E} \uput[ul](0.2,3.7){F} \uput[dr](4.5,2.7){G} \uput[ur](6.9,3.7){H}
\uput[d](2.6,2.2){O} \uput[ur](1.8,3.33){R} \uput[d](2.7,1.5){S} \uput[u](1.4,4.2){T}
\uput[ul](2.2,2){$\vect{\imath}$}\uput[ur](3.4,2.03){$\vect{\jmath}$}\uput[l](2.6,2.82){$\vect{k}$}
\end{pspicture}
\end{center}

Dans ce repère on a, en particulier C(6~;~8~;~0), F(6~;~0~;~4) et G(6~;~8~;~4).

Une des œuvres exposées est un triangle de verre représenté par le triangle ART qui a pour

sommets A(6~;~0~;~2), R(6~;~3~;~4) et T(3~;~0~;~4), Enfin, S est le point de coordonnées $\left(3~;~\dfrac52~;~0\right)$.

\begin{enumerate}
\item 
	\begin{enumerate}
		\item Vérifier que le triangle ART est isocèle en A.
		\item Calculer le produit scalaire $\vect{\text{AR}} \cdot \vect{\text{AT}}$.
		\item En déduire une valeur approchée à $0,1$ degré près de l'angle $\widehat{\text{RAT}}$.
	\end{enumerate}
\item 
	\begin{enumerate}
		\item Justifier que le vecteur $\vect{n}\begin{pmatrix}2\\-2\\3\end{pmatrix}$ est un vecteur normal au plan (ART).
		\item En déduire une équation cartésienne du plan (ART).
	\end{enumerate}
\item Un rayon laser dirigé vers le triangle ART est émis du plancher à partir du point S. On admet que ce rayon est orthogonal au plan (ART).
	\begin{enumerate}
		\item Soit $\Delta$ la droite orthogonale au plan (ART) et passant par le point S.
		
Justifier que le système ci-dessous est une représentation paramétrique de la droite $\Delta$ :

\[\left\{\begin{array}{l c r}
x&=&3+2k\\
y&=& \dfrac52 - 2k\\
z &=& 3k
\end{array}\right.,\: \text{avec }\:k \in \R.\]

		\item Soit L le point d'intersection de la droite $\Delta$, avec le plan (ART).

Démontrer que L a pour coordonnées $\left(5~;~\dfrac12~;~3\right)$.
	\end{enumerate}
\item L'artiste installe un rail représenté  par le segment  [DK] ou K est le milieu du segment [EH].

Sur ce rail, il positionne une source lumineuse laser en un point N du segment [DK] et il oriente ce
second rayon laser vers le point S.

\begin{center}
\psset{unit=1cm,arrowsize=2pt 3}
\begin{pspicture}(7.2,5.3)

\pspolygon(0.2,1.2)(4.5,0.2)(4.5,2.7)(0.2,3.7)
\psline(4.5,0.2)(6.9,1.2)(6.9,3.7)(4.5,2.7)
\psline(6.9,3.7)(2.6,4.7)(0.2,3.7)
\psline(0.2,2.45)(1.8,3.33)(1.4,4.2)
\psline[linestyle=dashed,linewidth=1.5pt](0.2,2.45)(1.4,4.2)
\psline[linestyle=dotted,linewidth=1.5pt](0.2,1.2)(2.6,2.2)(2.6,4.7)
\psline[linestyle=dotted,linewidth=1.5pt](2.6,2.2)(6.9,1.2)
\psline[linestyle=dotted,linewidth=1.5pt](4.75,4.2)(6.9,1.2)
\psline{->}(2.6,2.2)(2.2,2.)
\psline{->}(2.6,2.2)(3.4,2.03)
\psline{->}(2.6,2.2)(2.6,2.82)
\psline[ArrowInside=->](5.2,3.6)(2.7,1.5)(0.8,3.1)
\psdots(5.2,3.6)(2.7,1.5)(0.8,3.1)
\uput[l](0.2,2.45){A} \uput[dl](0.2,1.2){B} \uput[d](4.5,0.2){C} \uput[dr](6.9,1.2){D}
\uput[u](2.6,4.7){E} \uput[ul](0.2,3.7){F} \uput[dr](4.5,2.7){G} \uput[ur](6.9,3.7){H}
\uput[d](2.6,2.2){O} \uput[ur](1.8,3.33){R} \uput[d](2.7,1.5){S} \uput[u](1.4,4.2){T}
\uput[u](2.2,2){\small $\vect{\imath}$}\uput[u](3.4,2.03){\small $\vect{\jmath}$}\uput[u](2.6,2.82){\small $\vect{k}$}
\uput[ur](5.2,3.6){N}\uput[u](4.75,4.2){K}\uput[d](0.8,3.1){L}
\end{pspicture}
\end{center}

	\begin{enumerate}
		\item Montrer que, pour tout réel $t$ de l'intervalle [0~;~1], le point N de coordonnées $(0~;~8 - 4t~;~4t)$ est un point du segment [DK].
		\item  Calculer les coordonnées exactes du point N tel que les deux rayons laser représentés par les segments [SL] et [SN] soient perpendiculaires.
	\end{enumerate}
\end{enumerate}

\bigskip
